\section*{Uso de IA Generativa}

% TODO: manter apenas a declaração verdadeira para a versão submetida.
% O edital exige que qualquer uso de inteligência artificial generativa seja declarado.
% Opção 1: declarar que nenhuma ferramenta foi utilizada.
% Opção 2: informar ferramenta, partes da proposta e finalidade de uso.
\textbf{[TODO: inserir a declaração de uso de IA generativa aplicável.]}

\section{Introdução}

A \textit{Araucaria angustifolia}, conhecida como araucária ou pinheiro-do-paraná, é uma espécie característica da Floresta Ombrófila Mista, formação florestal associada principalmente à região Sul do Brasil \cite{saad2021}. Na Lista Nacional Oficial de Espécies da Flora Ameaçadas de Extinção, a espécie está classificada como Em Perigo (EN), conforme o Anexo I da Portaria GM/MMA nº 300, de 13 de dezembro de 2022 \cite{mma2022}.

A coleta do pinhão é sazonal e constitui fonte de renda complementar para agricultores familiares e comunidades rurais \cite{ipa2024}. A atividade envolve localizar as árvores e reconhecer as pinhas maduras, cuja produção varia entre indivíduos e ciclos reprodutivos. Demétrio et al. (2022) estudaram a produção de estróbilos masculinos e observaram essas estruturas em imagens de drone por avaliação manual, sem classificar automaticamente os indivíduos \cite{demetrio2022}.

Produzir um inventário georreferenciado dos indivíduos exige superar dificuldades de detecção e localização. Em floresta, o dossel pode ocultar árvores e reduzir a precisão de posições obtidas por GPS; é necessário também evitar omissões e duplicidades ao associar cada copa a um registro. Venancio, Disperati e Neves (2010) avaliaram o posicionamento de araucárias com receptores GPS \cite{venancio2010}. \citeauthor{karasinski2025} (\citeyear{karasinski2025}) detectaram araucárias em imagens do Google Earth e produziram mapas georreferenciados. Esses estudos oferecem referências para localização e detecção, mas não integram a classificação reprodutiva individual à avaliação do erro posicional. Um inventário que associe cada árvore à sua localização e às evidências reprodutivas poderia futuramente apoiar coletores no planejamento e acompanhamento das áreas de coleta; sua utilidade dependeria da qualidade e da completude dos registros.

A identificação de árvores individuais por suas copas é investigada por métodos de sensoriamento remoto e aprendizado profundo. Zhao et al. (2023) revisaram métodos de detecção e delimitação de copas com redes neurais convolucionais e apontaram desafios ligados à diversidade dos ambientes e dos dados \cite{zhao2023}. Trabalhos anteriores demonstraram a viabilidade de identificar araucárias em imagens com modelos de aprendizado profundo. \citet{karasinski2025} obtiveram mAP de 90,79\% na detecção da espécie em imagens do Google Earth. Em um fragmento florestal analisado com imagens de UAV, \citet{correa2024} reportou precisão de 79,3\% e revocação de 86,8\%. Esses resultados foram obtidos em condições diferentes e não são diretamente comparáveis, mas indicam que o desempenho deve ser avaliado conforme as características das imagens e da área estudada.

Além da detecção, a construção de um inventário individual exige associar cada copa à sua localização e às evidências de sua condição reprodutiva. Isso demanda investigar se modelos ajustados com imagens da área de estudo podem melhorar a identificação das araucárias e contribuir para uma classificação reprodutiva confiável. Nesse contexto, a pesquisa busca responder: como adaptar e avaliar modelos de aprendizado profundo para identificar e classificar araucárias por sexo em imagens de UAV, associando cada indivíduo à sua localização geográfica?

Para responder a essa pergunta, propõe-se como objetivo geral desenvolver e avaliar, em uma área piloto de Floresta Ombrófila Mista, um método computacional para produzir um inventário georreferenciado de araucárias individuais, com identificação das copas e classificação reprodutiva baseada em evidências visuais. Como objetivos específicos, pretende-se: (i) ajustar modelos de aprendizado profundo pré-treinados com imagens anotadas de araucárias para detectar e segmentar suas copas; (ii) ajustar e avaliar modelos pré-treinados com imagens anotadas de estruturas reprodutivas, obtidas em voos de aproximação com um drone menor, para classificar os indivíduos como masculinos, femininos ou indeterminados; (iii) associar imagens e evidências reprodutivas à copa correspondente; (iv) gerar registros geográficos individuais; e (v) avaliar o desempenho da detecção e classificação, a completude do inventário e o erro de localização, utilizando referências de campo conforme sua disponibilidade.

\section{Fundamentação teórica / trabalhos relacionados}

Estudos anteriores investigaram a detecção de araucárias por imagens do Google Earth e de UAV, empregando detectores da família YOLO, além de abordagens com imagens de satélite e dados LiDAR \cite{karasinski2025,correa2024,saad2021}. Esses trabalhos mostram possibilidades de mapeamento em diferentes ambientes, mas seus resultados não são diretamente comparáveis quando variam a área, a resolução das imagens, os rótulos e os protocolos de validação.

Também há pesquisas que utilizam imagens obtidas por drone para observar estruturas reprodutivas da espécie \cite{demetrio2022}. A classificação automática dessas estruturas e sua associação à árvore correta são problemas distintos da detecção da copa. Além disso, a ausência de estruturas visíveis em uma imagem não permite concluir a condição reprodutiva do indivíduo. Por isso, o inventário deverá admitir a classe indeterminada e registrar a imagem e a data que sustentam cada classificação.

A contribuição investigada é a integração entre identificação de copas, classificação reprodutiva baseada em evidência visual e geração de registros georreferenciados com erro de localização avaliado. A avaliação será realizada no mesmo contexto e com referências comuns, sem comparar diretamente métricas publicadas sob protocolos distintos.

\section{Metodologia}

O estudo será conduzido em uma área piloto de Floresta Ombrófila Mista, cuja disponibilidade ainda deverá ser confirmada. Serão utilizadas imagens aéreas ou de drone para delimitar copas e identificar indivíduos de \textit{Araucaria angustifolia}. Quando disponíveis, dados de altura da copa e imagens de outros ângulos poderão complementar a informação RGB. O desempenho será comparado com uma linha de base baseada em detector de objetos, como YOLO, usando as mesmas áreas e referências.

Em uma segunda etapa, imagens próximas das copas identificadas serão analisadas para reconhecer estróbilos masculinos ou pinhas femininas. A associação entre cada imagem próxima e a copa correspondente será verificada. Evidência visível reconhecida e atribuída com segurança poderá sustentar uma classificação masculina ou feminina; casos sem evidência suficiente ou com sinais contraditórios permanecerão indeterminados ou serão encaminhados para revisão humana.

As posições serão obtidas do levantamento georreferenciado e comparadas, em uma amostra, com medições de campo para avaliar o erro de localização. Os dados serão organizados em registros individuais contendo identificador, coordenadas, espécie identificada, classe reprodutiva atribuída, imagem de evidência e data da observação. Informações sobre a precisão posicional serão registradas conforme os dados disponíveis. O produto será exportável em formatos geográficos abertos, como GeoJSON e KML.

A avaliação considerará precisão, revocação, F1 e omissões na detecção; métricas por classe e taxa de indeterminados na classificação reprodutiva; erro de localização em metros; completude do inventário; e tempo necessário por área. A separação entre treino e teste será feita por árvore ou por área, evitando que imagens do mesmo indivíduo apareçam nos dois conjuntos. Árvores encobertas pelo dossel serão registradas como limite do método com imagens ópticas aéreas.

\section{Resultados esperados}

Espera-se obter um inventário georreferenciado de araucárias individuais e um fluxo reproduzível para produzi-lo, acompanhados de uma avaliação dos erros de detecção, classificação reprodutiva e localização. O inventário deverá preservar a evidência visual e a data de cada observação, permitir casos indeterminados e ser exportável em formatos geográficos abertos.

O estudo deverá indicar em que condições a integração de imagens de copa e de estruturas reprodutivas é viável, quais erros permanecem e como a cobertura do dossel afeta a completude. A disponibilidade de área, imagens adequadas e referências de campo será determinante para o escopo final. O artefato será apresentado como ferramenta experimental de inventário e apoio à pesquisa, não como levantamento validado para qualquer área sem avaliação local.
